\documentclass[10pt]{article}

% --- Engine-Adaptive Preamble (Supports pdfLaTeX in Overleaf & LuaLaTeX) ---
\usepackage[letterpaper, top=1in, bottom=1in, left=1.5in, right=1.5in, textwidth=5.5in, textheight=9in]{geometry}
\usepackage{iftex}
\ifPDFTeX
  \usepackage[utf8]{inputenc}
  \usepackage[T1]{fontenc}
  \usepackage{times}
\else
  \usepackage{fontspec}
  \usepackage[english, bidi=basic, provide=*]{babel}
  \babelprovide[import, onchar=ids fonts]{english}
  \babelfont{rm}{Noto Serif}
  \babelfont{sf}{Noto Sans}
  \babelfont{tt}{Noto Sans Mono}
\fi

\usepackage{booktabs}       % professional-quality tables
\usepackage{amsfonts}       % blackboard math symbols
\usepackage{amsmath}        % math equations and operators
\usepackage{nicefrac}       % compact symbols for 1/2, etc.
\usepackage{microtype}      % microtypography
\usepackage{xcolor}         % colors
\usepackage{tabularx}       % auto-wrapping columns within \linewidth
\usepackage{adjustbox}      % shrink wide tables cleanly to max width=\linewidth
\usepackage{graphicx}       % figures
\usepackage[numbers, compress]{natbib}
\usepackage{titlesec}
\usepackage{parskip}
\usepackage{caption}
\usepackage{enumitem}
\usepackage{hyperref}       % hyperlinks (always last)

\graphicspath{{figures/}}

% --- Replicate NeurIPS 2026 Visual Layout & Typography ---
\setlength{\parindent}{0pt}
\setlength{\parskip}{5.5pt}
\linespread{1.05} % 10pt type with ~11pt leading

\titleformat{\section}{\normalfont\fontsize{12}{14}\bfseries}{\thesection}{1em}{}
\titleformat{\subsection}{\normalfont\fontsize{10}{12}\bfseries}{\thesubsection}{1em}{}
\titleformat{\subsubsection}{\normalfont\fontsize{10}{12}\bfseries}{\thesubsubsection}{1em}{}

% Table captions: 1 line space before title, 1 line space after title
\captionsetup[table]{position=top, skip=11pt, aboveskip=11pt, belowskip=11pt, font=small}
% Figure captions: below the figure, same size as table captions
\captionsetup[figure]{position=bottom, skip=6pt, font=small}

% Checklist helper macros matching NeurIPS 2026 style
\newcommand{\answerYes}[1][]{\textcolor{blue}{[Yes]} #1}
\newcommand{\answerNo}[1][]{\textcolor{orange}{[No]} #1}
\newcommand{\answerNA}[1][]{\textcolor{gray}{[N/A]} #1}

\makeatletter
\renewcommand{\maketitle}{%
  \vbox{%
    \hsize\textwidth
    \linewidth\hsize
    \vskip 0.1in
    \hrule height 4pt
    \vskip 0.25in
    \centering
    {\fontsize{17}{20}\selectfont\bfseries \@title \par}
    \vskip 0.25in
    \hrule height 1pt
    \vskip 0.25in
    {\normalsize \@author \par}
    \vskip 0.3in
  }%
}
% Abstract: Two line spaces precede abstract; indented 0.5in left & right; 12pt bold centered title
\renewenvironment{abstract}{%
  \vskip 22pt%
  \centerline{\fontsize{12}{14}\selectfont\bfseries Abstract}%
  \vspace{11pt}%
  \begin{list}{}{%
    \setlength{\leftmargin}{0.5in}%
    \setlength{\rightmargin}{0.5in}%
    \setlength{\topsep}{0pt}%
    \setlength{\partopsep}{0pt}%
  }\item\relax\fontsize{10}{11}\selectfont
}{%
  \end{list}%
  \vskip 11pt%
}
\newenvironment{ack}{%
  \section*{Acknowledgments and disclosure of funding}%
}{}
\makeatother

\DeclareMathOperator*{\argmax}{arg\,max}
\DeclareMathOperator{\norm}{norm}
\DeclareMathOperator{\TV}{TV}

\title{JCSD: Hardware-Aware Joint Cascade and Speculative-Tree Decoding for Fast LLM Inference}

\author{%
  \begin{tabular}{ccc}
    \textbf{Ayanna Mondal} &
    \textbf{Swastik Pramanik} &
    \textbf{Takshay Bansal}\thanks{HDL Course Project, IIT Guwahati. Code: \url{https://github.com/TakshayBansal/hdl_project} (branch \texttt{complete-pipeline}).} \\
    Roll No.: 230102015 &
    Roll No.: 230102094 &
    Roll No.: 230102096 \\
    \texttt{m.ayanna@iitg.ac.in} &
    \texttt{p.swastik@iitg.ac.in} &
    \texttt{t.bansal@iitg.ac.in} \\[1.2ex]
    \multicolumn{3}{c}{%
      \begin{tabular}{cc}
        \textbf{Youkta Mandavkar} &
        \textbf{Shirshendu Sengupta} \\
        Roll No.: 230108062 &
        Roll No.: 230108072 \\
        \texttt{m.youkta@iitg.ac.in} &
        \texttt{g.shirshendu@iitg.ac.in}
      \end{tabular}%
    }\\[1.5ex]
    \multicolumn{3}{c}{\normalfont Indian Institute of Technology Guwahati, Guwahati, Assam, India}
  \end{tabular}
}

\begin{document}

\maketitle

\begin{abstract}
Autoregressive decoding in large language models (LLMs) is severely memory-bandwidth bound, leaving GPU compute underutilized during token-by-token generation. While speculative token trees (SEQUOIA) accelerate exact sampling via dynamic-programming-optimized tree topologies, and speculative cascades (Faster Cascades) trade small, controlled quality tolerances for higher draft acceptance via a deferral parameter $\alpha$, prior work optimizes these dimensions independently. However, $\alpha$ directly governs token acceptance probabilities, which dictate the optimal tree geometry under specific hardware latency curves. We propose \textbf{JCSD} (\emph{Joint Cascade + Speculative-Tree Decoding}), a hardware-aware framework that jointly optimizes the cascade deferral parameter $\alpha$ and token-tree topology $T$ under an explicit quality constraint. By exploiting a structural decoupling property—output quality depends solely on $\alpha$ rather than tree shape—we decompose a costly three-way search into a single offline quality/acceptance pass, lightweight per-GPU latency profiling, and a fast dynamic-programming search. Evaluated on Pythia-160M $\to$ Pythia-2.8B on CNN/DailyMail using an NVIDIA Tesla T4 (FP16), JCSD achieves a \textbf{1.48$\times$} wall-clock speedup over autoregressive decoding, outperforming sequential tuning (\textbf{1.41$\times$}, paired latency reduction of \textbf{1.41 ms/token}, 95\% CI $[1.02, 1.79]$), SEQUOIA (\textbf{1.37$\times$}), and Faster Cascades chains (\textbf{1.35$\times$}). On a second platform (the same T4 in FP32, where verifying 64 tokens costs $3.5\times$ one token instead of $1.4\times$), the optimal trees change sharply (SEQUOIA: 10 $\to$ 2 nodes; joint: depth 3 $\to$ 5), transplanting the FP16 joint optimum loses \textbf{5.9\%} measured speed, and joint search beats sequential by \textbf{16\%} ($1.375\times$ vs.\ $1.185\times$). On FP32, however, a Faster Cascades chain, which lies inside the joint search space, measured 2.9\% faster than the predicted joint optimum; we trace this ranking error to a draft pass that the cost model does not price.
\end{abstract}

\section{Introduction}

Generating text with large language models (LLMs) requires loading the full model weights from GPU high-bandwidth memory (HBM) for every single generated token. For a 2.8B-parameter model in FP16 ($\approx 5.6$\,GB), each forward pass at batch size 1 is heavily memory-bandwidth bound, leaving tensor cores largely idle. Our profiling on an NVIDIA Tesla T4 confirms this bottleneck: verifying 1 token with Pythia-2.8B takes $30.1$\,ms, whereas verifying 64 tokens simultaneously takes only $42.7$\,ms—a $1.4\times$ latency increase for $64\times$ the compute.

Two prominent research lines exploit this excess compute capacity:
\begin{itemize}
    \item \textbf{Speculative decoding with token trees (SEQUOIA)} \citep{chen2024sequoia}: A lightweight draft model $q$ proposes a tree of candidate continuations, which the target model $p$ verifies in a single forward pass using a tree attention mask. Sampling without replacement and dynamic programming (DP) over hardware latency curves yield exact samples from $p$ at substantially higher throughput.
    \item \textbf{Speculative cascades (Faster Cascades)} \citep{narasimhan2025faster}: Rather than enforcing exact equivalence to $p$, verification targets a blended distribution $\pi_\alpha$ governed by a deferral knob $\alpha \in [0, 1]$. Relaxing strict target matching increases draft acceptance rates while maintaining controlled output quality.
\end{itemize}

Prior work treats tree topology and cascade deferral in isolation (Table~\ref{tab:gap}). Yet these choices are tightly coupled: changing $\alpha$ shifts the positional acceptance probabilities $p_k(\alpha)$, shifting $p_k(\alpha)$ alters the optimal branching factor and depth of the draft tree, and the target hardware $H$ prices the marginal cost of each draft level and verified node.

\begin{table}[h]
  \caption{Comparison of speculative decoding paradigms. JCSD is the first framework to jointly optimize cascade deferral $\alpha$ and tree topology $T$ under a hardware cost model and quality constraint.}
  \label{tab:gap}
  \centering
  \small
  \begin{tabularx}{\linewidth}{@{}lXX@{}}
    \toprule
    Method & Optimized variables & Unmodeled / fixed aspects \\
    \midrule
    Faster Cascades \citep{narasimhan2025faster} & Deferral rule \& $\alpha$ (quality vs.\ acceptance) & Fixed chain topology; ignores hardware cost curves \\
    SEQUOIA \citep{chen2024sequoia} & Tree shape (DP) \& size (hardware profile) & No quality trade-off; target fixed strictly to $p$ \\
    \textbf{JCSD (ours)} & \textbf{Joint $(\alpha, T)$ under quality budget, per $H$} & \textbf{---} \\
    \bottomrule
  \end{tabularx}
  \vspace{11pt}
\end{table}

\paragraph{Contributions.} In this report, we present the design, implementation, verification, and empirical evaluation of JCSD:
\begin{enumerate}
    \item \textbf{Constrained joint formulation and decoupling:} We formulate hardware-aware $(\alpha, T)$ optimization under a quality constraint and identify a decoupling property that separates GPU quality/acceptance measurement from CPU tree search.
    \item \textbf{Unified, verified inference engine:} We implement a batched TokenV3 verifier, DP tree optimizer, KV-cache surgery engine, and a shared tree decoder that runs all baselines through identical code paths, validated by 84 unit and statistical tests.
    \item \textbf{Empirical validation on two hardware platforms:} We evaluate four core hypotheses (H1--H4) on Pythia-160M $\to$ Pythia-2.8B (CNN/DailyMail) on a T4 in FP16 and in FP32, showing that joint optimization outperforms sequential tuning in measured wall-clock speed on both ($1.481\times$ vs.\ $1.405\times$; $1.375\times$ vs.\ $1.185\times$) and that the optimum does not transfer across platforms without loss.
    \item \textbf{Diagnostic analysis of modeling assumptions:} We validate and calibrate the analytical cost model against end-to-end timings, identify an unmodeled draft pass that explains its largest errors, quantify sibling waste in static trees (63--80\%), and measure depth-dependent drift in positional acceptance.
\end{enumerate}

\section{Background and related work}

\paragraph{Speculative decoding.} Introduced by \citet{leviathan2023fast} and \citet{chen2023accelerating}, speculative decoding uses a fast draft model $q$ to generate candidate tokens and a larger target model $p$ to score them in parallel. A draft token $x \sim q$ is accepted with probability $\min(1, p(x)/q(x))$; upon the first rejection, a replacement token is drawn from the residual distribution $\norm(\max(p - q, 0))$. This guarantees that the resulting sequence follows $p$ identically.

\paragraph{Token trees (SEQUOIA).} Sequential chains waste computation whenever an early draft token is rejected. SEQUOIA \citep{chen2024sequoia} generalizes chains to trees, sampling $k$ distinct children per node without replacement from $q$ and verifying all nodes in a single target pass via custom tree attention masks. Under a positional acceptance assumption, SEQUOIA uses dynamic programming to find the expected-token-maximizing tree of size $n$ and selects $n$ from measured hardware latencies.

\paragraph{Speculative cascades (Faster Cascades).} Cascades route easy tokens to a small model and hard tokens to a large model. \citet{narasimhan2025faster} embed deferral directly inside speculative verification by defining a composite target distribution $\pi_\alpha$ that blends $q$ and $p$ based on target confidence. Because speculative sampling is exact for any target distribution, verifying against $\pi_\alpha$ produces exact samples from $\pi_\alpha$ with higher acceptance rates than $p$.

\section{Problem formulation}

\subsection{Constrained hardware-aware objective}
An unconstrained maximization of expected tokens per unit time $\mathbb{E}[\text{tokens}] / \text{Cost}$ degenerates to $\alpha \to 1$, where every draft token is accepted and output quality collapses to that of the small draft model. Furthermore, hardware $H$ is a fixed deployment environment rather than a decision variable. For a given hardware platform $H$, we formulate JCSD as:
\begin{equation}
  \max_{\alpha \in [0, 1],\, T} \; \mathrm{Speedup}(\alpha, T; H) \quad \text{subject to} \quad Q(\alpha) \ge (1 - \mathrm{tol})\, Q(p),
  \label{eq:obj}
\end{equation}
where $Q(\alpha)$ is the task quality metric under target $\pi_\alpha$, $Q(p)$ is the exact target model quality, and $\mathrm{tol} \in \{2\%, 5\%, 10\%\}$ specifies the maximum allowable relative quality degradation.

\subsection{TokenV3 cascade target distribution}
Following \citet{narasimhan2025faster}, let $p$ denote raw (temperature-1) target probabilities used for deferral gating, and let $q_s, p_s$ denote temperature-scaled draft and target distributions. For deferral parameter $\alpha \in [0, 1]$, we define the plausible token set $\mathcal{T}_\alpha$ and cascade target $\pi_\alpha$:
\begin{align}
  \mathcal{T}_\alpha &= \left\{\, v \in \mathcal{V} : p(v) \ge (1 - \alpha)\,\max_{v' \in \mathcal{V}} p(v') \,\right\}, \\
  \pi_\alpha(v) &= q_s(v)\,\mathbb{1}[v \in \mathcal{T}_\alpha] + p_s(v) \sum_{u \notin \mathcal{T}_\alpha} q_s(u).
\end{align}
Intuitively, $\pi_\alpha$ retains the draft model's probability mass on tokens that the target model deems within a $(1-\alpha)$ factor of its top choice, while reallocating implausible draft mass to $p_s$. Setting $\alpha = 1$ recovers $\pi_1 = q_s$. Note that $\alpha = 0$ restricts $\mathcal{T}_0$ to $\arg\max p$ rather than recovering exact $p$; therefore, exact target decoding ($p$) is evaluated as a distinct baseline setting.

\subsection{Tree verification, expected yield, and latency cost}
At any tree node, $k$ children $c_1, \dots, c_k$ are sampled sequentially without replacement from $q$ via the Gumbel top-$k$ trick. Initializing $R \leftarrow \pi_\alpha$ and $D \leftarrow q$, for $i = 1, \dots, k$, $D$ is renormalized and child $c_i$ is accepted with probability $\min(1, R(c_i)/D(c_i))$. On rejection, $R \leftarrow \norm(\max(R - D, 0))$ and $D(c_i) \leftarrow 0$; if all $k$ children reject, a token is sampled from the final residual $R$. This procedure produces an exact sample from $\pi_\alpha$.

Let $p_k(\alpha)$ denote the probability that the $k$-th child of a visited node is accepted. Under SEQUOIA's positional assumption, the expected number of accepted tokens per verification step for tree $T$ is:
\begin{equation}
  G_\alpha(T) = \sum_{v \in T} \prod_{\substack{u \in \mathrm{path}(v) \\ u \neq \mathrm{root}}} p_{\mathrm{rank}(u)}(\alpha).
\end{equation}
For a tree with level widths $w_0 = 1, w_1, \dots, w_d$ and total nodes $n = \sum_{l=0}^d w_l$, we model wall-clock step latency and speedup relative to autoregressive step time $t_{\mathrm{AR}}$ as:
\begin{equation}
  \mathrm{Cost}(T; H) = t_{\mathrm{verify}}(n) + \sum_{l=0}^{d-1} t_{\mathrm{draft}}(w_l) + c_0 + c_1\, G_\alpha(T), \qquad
  \mathrm{Speedup}(\alpha, T; H) = \frac{G_\alpha(T)\, t_{\mathrm{AR}}}{\mathrm{Cost}(T; H)},
  \label{eq:cost}
\end{equation}
where $t_{\mathrm{verify}}$ and $t_{\mathrm{draft}}$ are measured hardware latencies, level $0$ is priced at width $\lceil G_\alpha(T) \rceil$ (accounting for drafting from the chunk of tokens committed in the prior step), the deepest level $d$ requires no draft pass, and $c_0, c_1 \ge 0$ capture runtime orchestration overhead fitted via non-negative least squares.

\subsection{Joint vs.\ sequential optimization and hypotheses}
Let $\mathcal{A} = \{\alpha : Q(\alpha) \ge (1 - \mathrm{tol})\, Q(p)\}$ be the feasible set. We compare:
\begin{align}
  \text{Joint:} &\quad \max_{\alpha \in \mathcal{A}} \max_{T} \mathrm{Speedup}(\alpha, T; H), \\
  \text{Sequential:} &\quad T_p^\star = \argmax_{T} \mathrm{Speedup}_p(T), \quad \max_{\alpha \in \mathcal{A}} \mathrm{Speedup}(\alpha, T_p^\star; H).
\end{align}
We test four formal hypotheses:
\begin{itemize}
    \item \textbf{H1:} Acceptance probabilities $P(\alpha)$ vary systematically with $\alpha$, with first-child acceptance $p_1(\alpha)$ increasing monotonically with $\alpha$.
    \item \textbf{H2:} Optimal tree topology $T^\star(\alpha)$ shifts with $\alpha$, becoming deeper and narrower at the top as $p_1(\alpha)$ rises.
    \item \textbf{H3:} Optimal $(\alpha^\star, T^\star)$ depends on hardware platform $H$, incurring a speedup penalty when transferring configurations across platforms.
    \item \textbf{H4:} Joint search outperforms sequential search in \emph{measured} wall-clock speedup at equal output quality.
\end{itemize}

\section{Method and system architecture}

\subsection{The decoupling insight}
Exhaustively searching over $(\alpha, T, H)$ via end-to-end tree decoding on GPUs would be computationally prohibitive. However, because tree verification is mathematically exact for any target distribution $\pi_\alpha$:
\begin{enumerate}
    \item \textbf{Output quality depends only on $\alpha$}, independent of tree topology $T$ and hardware $H$.
    \item \textbf{Quality $Q(\alpha)$ and positional acceptance $P(\alpha)$} can be measured simultaneously via batched, token-by-token sampling from $\pi_\alpha$—without constructing trees—\emph{once} per (model pair, task, temperature) and reused across all hardware platforms.
    \item \textbf{Joint optimization} reduces to lightweight per-GPU latency profiling ($t_{\mathrm{verify}}(n), t_{\mathrm{draft}}(n)$) followed by an offline CPU dynamic program (capped at $N=128$ nodes, depth $D=12$, branching factor $B=16$) and table lookup taking seconds.
\end{enumerate}

\subsection{End-to-end pipeline and decoder invariants}
Our pipeline consists of seven modular stages: (0) \texttt{check\_tree\_attention.py} verifies custom tree mask support on real models; (2) \texttt{run\_profile.py} measures CUDA event median latencies over 30 repetitions; (3--4) \texttt{run\_quality.py} performs left-padded batched sampling with inline float64 batched verification (\texttt{AcceptCounter}) to extract $Q(\alpha)$ and $P(\alpha)$; (5) \texttt{run\_search.py} runs the SEQUOIA DP and selects joint and sequential optima; (6) \texttt{run\_e2e.py} benchmarks wall-clock inference; (6b) \texttt{calibrate.py} fits overhead $(c_0, c_1)$ and depth-dependent acceptance matrices; and (7) \texttt{make\_figures.py} / \texttt{summarize.py} generate plots, \LaTeX{} tables, and statistical verdicts.

\paragraph{Shared decoder architecture.} Every speculative method—including chains (depth-$D$, width-1 trees), SEQUOIA (mode $p$), sequential, and joint—executes through a single unified \texttt{TreeDecoder} class. This eliminates implementation bias so measured latency differences reflect algorithmic choices alone. Key KV-cache invariants ensure exactness:
\begin{itemize}
    \item The target KV cache stores all committed tokens except the root; each verification step processes $[\text{root}] + \text{drafted nodes}$ with explicit \texttt{position\_ids} ($\text{ctx} + \text{depth}$) and ancestor-only tree attention masks.
    \item After verification, \texttt{keep\_path} performs in-place KV-cache surgery, compacting the accepted branch into contiguous prefix slots via a version-portable \texttt{crop\_to} helper.
    \item When a walk in cascade mode ends at a leaf on the deepest tree level, the draft distribution $q$ at that leaf is unknown (the deepest level is never drafted from), yet it is needed to form $\pi_\alpha$ for the bonus token sampled there. The accepted path is therefore fed through the draft model immediately. This keeps sampling exact but costs one extra draft forward pass, which the cost model of Eq.~\eqref{eq:cost} does not price (Section~\ref{sec:costmodel}).
\end{itemize}

\subsection{Verification strategy}
We enforced a strict test-driven discipline across 84 automated tests using synthetic GPT-NeoX models with sharpened output embeddings:
\begin{itemize}
    \item \textbf{Distribution exactness:} $\chi^2$ goodness-of-fit tests confirm that the batched float64 verifier ($30{,}000$ samples per configuration, including a draft with zero support) and the multi-step \texttt{TreeDecoder} ($4{,}000$ decodes per configuration, on both the first token and the second token through KV-cache surgery) sample from $\pi_\alpha$ exactly.
    \item \textbf{Tree mask equivalence:} Per-node logits under custom 4D tree attention masks match independent path forward passes on both synthetic models and real Pythia weights on GPU (context lengths 64 and 400).
    \item \textbf{Baseline identity:} A root-only tree produces token-identical outputs to autoregressive decoding under matched seeds.
\end{itemize}

\section{Experimental setup}

Table~\ref{tab:setup} summarizes our experimental configuration. To prevent thermal throttling or GPU clock drift from biasing wall-clock comparisons, \texttt{run\_e2e.py} interleaves all methods on every prompt with rotated execution order and common random numbers across paired passes, excluding prompt prefill from per-token latency.

\begin{table}[h]
  \caption{Experimental setup and hyperparameters for JCSD evaluation on CNN/DailyMail.}
  \label{tab:setup}
  \centering
  \small
  \begin{tabularx}{\linewidth}{@{}lX@{}}
    \toprule
    Component & Specification \\
    \midrule
    Model pair A & Draft: \texttt{EleutherAI/pythia-160m} (12 layers); Target: \texttt{EleutherAI/pythia-2.8b} (32 layers) \\
    Task & CNN/DailyMail summarization (300-word article truncation, 64 new tokens, ROUGE-2) \\
    Dataset splits & Calibration (validation): 200 articles; End-to-end test: 100 articles; Warm-up: 2 train articles \\
    $\alpha$ grid & $\{0, 0.1, 0.2, 0.3, 0.5, 0.7, 0.9, 0.92, 0.95, 0.97, 0.99, 0.995, 0.998, 0.999, 1.0\}$ ($T = 1.0$) \\
    Quality budgets & Tolerances $\mathrm{tol} \in \{2\%, 5\%, 10\%\}$ relative to $Q(p)$ (mean criterion; paired CI sensitivity) \\
    DP \& simulation caps & $N = 128$ nodes, $D = 12$ depth, $B = 16$ children; 4 acceptance draws/pos ($\approx 50{,}000$/setting) \\
    Platform 1 & NVIDIA Tesla T4 (Google Colab), FP16, PyTorch 2.14.1+cu130, Transformers 5.22.1, SDPA \\
    Platform 2 & NVIDIA Tesla T4 (Kaggle), FP32, PyTorch 2.11.0+cu130, Transformers 5.20.2, SDPA (tensor-core vs.\ CUDA-core regime) \\
    E2E protocol & Batch size 1, 100 test prompts $\times$ 3 passes, interleaved \& order-rotated, common seeds \\
    \bottomrule
  \end{tabularx}
  \vspace{11pt}
\end{table}

Platform 2 is the same GPU model at a different numerical precision. FP16 runs matrix products on the T4's tensor cores ($\approx 65$\,TFLOPS), whereas FP32 uses its CUDA cores ($\approx 8$\,TFLOPS) and moves twice the bytes, which changes the shape of $t_{\mathrm{verify}}(n)$ (Section~\ref{sec:profile}). Quality and acceptance were measured once (in FP16) and reused for FP32, per the decoupling property.

\section{Results}

\subsection{Hardware latency profiles}
\label{sec:profile}
Table~\ref{tab:profile} and Figure~\ref{fig:profiles} report median forward latencies over 30 runs at context length 512 using the masked tree-attention path. Three hardware characteristics govern tree selection:
\begin{enumerate}
    \item \textbf{Wide verification sweet spot on FP16:} Verifying 64 tokens with Pythia-2.8B costs $42.7$\,ms versus $30.1$\,ms for 1 token ($1.42\times$), only hitting compute-bound scaling at $n \ge 128$.
    \item \textbf{Steep verification curve on FP32:} The same ratio is $3.5\times$ ($181.3$ vs.\ $51.6$\,ms), with a first step already at $n = 4$ ($71.0$\,ms). Wide trees are therefore expensive on Platform 2.
    \item \textbf{Launch-bound draft model:} Pythia-160M exhibits a nearly flat latency of $\approx 10$--$13$\,ms across $n \in [1, 128]$ on both platforms. Because each tree level requires a sequential draft pass costing one-third (FP16) to one-fifth (FP32) of a target verification pass, deep trees incur a steep latency penalty.
\end{enumerate}

\begin{table}[h]
  \caption{Measured forward-pass latency (ms) on NVIDIA Tesla T4 (context 512, tree-mask path). Profiled AR step latency is $33.9$\,ms (FP16) and $53.6$\,ms (FP32); measured end-to-end AR loop latency is $36.53$ and $50.03$\,ms/token respectively.}
  \label{tab:profile}
  \centering
  \begin{adjustbox}{max width=\linewidth}
  \begin{tabular}{llrrrrrrrrr}
    \toprule
    Platform & Tokens per forward ($n$) & 1 & 2 & 4 & 8 & 16 & 32 & 64 & 128 & 256 \\
    \midrule
    T4 FP16 & $t_{\mathrm{verify}}$ (Pythia-2.8B, ms) & 30.1 & 31.1 & 31.5 & 33.3 & 36.4 & 38.3 & 42.7 & 57.7 & 96.3 \\
            & $t_{\mathrm{draft}}$ (Pythia-160M, ms) & 10.0 & 10.4 & 10.0 & 12.6 & 10.2 & 10.4 & 10.4 & 9.7 & 13.8 \\
    \midrule
    T4 FP32 & $t_{\mathrm{verify}}$ (Pythia-2.8B, ms) & 51.6 & 53.0 & 71.0 & 88.5 & 87.7 & 98.0 & 181.3 & 316.0 & 542.9 \\
            & $t_{\mathrm{draft}}$ (Pythia-160M, ms) & 10.6 & 12.2 & 13.0 & 12.4 & 10.9 & 11.1 & 11.4 & 14.9 & 29.6 \\
    \bottomrule
  \end{tabular}
  \end{adjustbox}
  \vspace{11pt}
\end{table}

\begin{figure}[t]
  \centering
  \includegraphics[width=\linewidth]{fig1_profiles.png}
  \caption{Hardware profiles. Left: target verification latency $t_{\mathrm{verify}}(n)$; right: draft latency per tree level. FP32 verification grows far faster with $n$ than FP16, while the draft model is launch-bound on both platforms.}
  \label{fig:profiles}
\end{figure}

\subsection{Quality and acceptance vs.\ $\alpha$ (Hypothesis H1)}
Table~\ref{tab:quality} and Figure~\ref{fig:quality} present ROUGE-2 quality and positional acceptance rates across the calibration split. Remarkably, for all $\alpha \le 0.99$, TokenV3 achieves \emph{higher} ROUGE-2 than sampling directly from the target model $p$ ($0.0578$--$0.0773$ vs.\ $0.0476$, with paired 95\% CIs strictly above zero). At temperature $1.0$, raw sampling from Pythia-2.8B occasionally selects low-probability tail tokens; TokenV3 prunes this tail similarly to nucleus sampling, improving summary quality. Consequently, the quality constraint only becomes binding very close to $\alpha = 1$, crossing $Q(p) = 0.0476$ between $\alpha = 0.995$ ($Q = 0.0480$) and $\alpha = 0.998$ ($Q = 0.0422$). Figures~\ref{fig:quality}--\ref{fig:trees} therefore plot $\alpha$ on a logarithmic axis in $1-\alpha$.

\begin{table}[h]
  \caption{Hardware-independent ROUGE-2 quality $Q(\alpha)$ and positional acceptance rates on 200 CNN/DailyMail calibration prompts. Simulated $p_1$ matches analytic $1 - \mathrm{TV}(\pi_\alpha, q)$ within $0.002$ across all settings.}
  \label{tab:quality}
  \centering
  \begin{adjustbox}{max width=\linewidth}
  \begin{tabular}{lccccc}
    \toprule
    Setting & ROUGE-2 $Q$ & 95\% CI & $Q - Q(p)$ paired 95\% CI & $p_1$ & $p_2$ \\
    \midrule
    Draft only ($q$) & 0.0150 & $[0.0121, 0.0183]$ & --- & --- & --- \\
    \textbf{Target only ($p$)} & \textbf{0.0476} & $[0.0395, 0.0563]$ & --- & \textbf{0.636} & 0.093 \\
    $\alpha = 0.0$ & 0.0715 & $[0.0605, 0.0831]$ & $[+0.015, +0.034]$ & 0.666 & 0.100 \\
    $\alpha = 0.5$ & 0.0773 & $[0.0658, 0.0895]$ & $[+0.020, +0.040]$ & 0.680 & 0.107 \\
    $\alpha = 0.9$ & 0.0675 & $[0.0576, 0.0781]$ & $[+0.012, +0.029]$ & 0.695 & 0.117 \\
    $\alpha = 0.95$ & 0.0618 & $[0.0523, 0.0718]$ & $[+0.007, +0.022]$ & 0.717 & 0.115 \\
    $\alpha = 0.99$ & 0.0578 & $[0.0494, 0.0669]$ & $[+0.001, +0.020]$ & 0.769 & 0.103 \\
    \textbf{$\alpha = 0.995$} & \textbf{0.0480} & --- & $[-0.009, +0.009]$ & \textbf{0.794} & --- \\
    $\alpha = 0.998$ & 0.0422 & --- & $[-0.014, +0.004]$ & 0.826 & --- \\
    $\alpha = 0.999$ & 0.0421 & --- & $[-0.015, +0.003]$ & 0.848 & --- \\
    $\alpha = 1.0$ ($= q$) & 0.0150 & $[0.0121, 0.0183]$ & $[-0.041, -0.025]$ & 1.000 & 0.000 \\
    \bottomrule
  \end{tabular}
  \end{adjustbox}
  \vspace{11pt}
\end{table}

\begin{figure}[t]
  \centering
  \includegraphics[width=0.6\linewidth]{fig2_quality.png}
  \caption{Quality $Q(\alpha)$ with 95\% bootstrap CI, against the target-only and draft-only references. The cascade exceeds the target model until $\alpha \approx 0.995$.}
  \label{fig:quality}
\end{figure}

\begin{figure}[t]
  \centering
  \includegraphics[width=\linewidth]{fig3_acceptance.png}
  \caption{Left: acceptance vector $p_k(\alpha)$ (H1); the dashed line is $p_1$ for the exact target. Right: per-position distribution of $p_1 = 1 - \mathrm{TV}(\pi_\alpha, q)$, which is bimodal (at $\alpha=0$, 23\% of positions exceed 0.95 and 27\% fall below 0.5).}
  \label{fig:acceptance}
\end{figure}

\paragraph{Verdict on H1 (supported).} First-child acceptance $p_1(\alpha)$ increases strictly monotonically with $\alpha$ (Spearman $\rho = 1.000$), rising from $0.666$ ($\alpha=0$) to $0.794$ ($\alpha=0.995$) and $0.848$ ($\alpha=0.999$), with most variation concentrated in the logarithmic regime $1 - \alpha \le 10^{-2}$ (Figure~\ref{fig:acceptance}). The average hides strong bimodality across contexts, which matters for the positional assumption (Section~\ref{sec:costmodel}).

\subsection{Optimal tree topology shift (Hypothesis H2)}
Table~\ref{tab:trees} and Figure~\ref{fig:trees} show how the DP-optimal tree structure evolves with $\alpha$ on both platforms. On FP16, for exact $p$ ($p_1 = 0.636$), high rejection risk favors a shallow, wide tree of depth $2$ and $10$ nodes ($[1, 4, 5]$). Once $\alpha \ge 0.99$ pushes $p_1 \ge 0.769$, deeper speculation becomes profitable: the optimal tree narrows its first branching factor from $4 \to 3$ and adds a third level ($[1, 3, 4, 4]$, $n=12, d=3$). On FP32, where wide verification is expensive, the optimum for exact $p$ is a single drafted token ($n=2$, $d=1$); at $\alpha = 0.995$ it jumps to a deep, slowly widening tree ($[1, 2, 3, 4, 5, 6]$, $n=21$, $d=5$). Figure~\ref{fig:surface} shows the predicted speedup surface behind these choices.

\begin{table}[h]
  \caption{DP-optimal token-tree topologies across $\alpha$ on both platforms, showing level widths, total nodes $n$, depth $d$, and expected accepted tokens per step $G_\alpha(T)$.}
  \label{tab:trees}
  \centering
  \begin{adjustbox}{max width=\linewidth}
  \begin{tabular}{lllccc}
    \toprule
    Platform & $\alpha$ setting & Optimal tree level widths & Nodes ($n$) & Depth ($d$) & Expected tokens $G$ \\
    \midrule
    T4 FP16 & Exact $p$ (SEQUOIA) & $[1, 4, 5]$ & 10 & 2 & 2.38 \\
            & $\alpha \in [0.0, 0.97]$ & $[1, 4, 5]$ or $[1, 4, 7]$ & 10--12 & 2 & 2.48--2.63 \\
            & $\alpha \in [0.99, 0.998]$ & $[1, 3, 4, 4]$ & 12 & 3 & 3.31--3.49 \\
            & $\alpha = 0.999$ & $[1, 3, 4, 4, 5, 6]$ & 23 & 5 & 4.91 \\
            & $\alpha = 1.0$ (draft only) & $[1, 1, \dots, 1]$ (13-node chain) & 13 & 12 & 13.00 \\
    \midrule
    T4 FP32 & Exact $p$ (SEQUOIA) & $[1, 1]$ & 2 & 1 & 1.64 \\
            & $\alpha \in [0.0, 0.97]$ & $[1, 1]$ & 2 & 1 & 1.67--1.72 \\
            & $\alpha = 0.99$ & $[1, 1, 1]$ & 3 & 2 & 2.36 \\
            & $\alpha \in [0.995, 0.998]$ & $[1, 2, 3, 4, 5, 6]$ & 21 & 5 & 4.45--4.68 \\
            & $\alpha = 0.999$ & $[1, 2, 3, 4, 5, 6, 7]$ & 28 & 6 & 5.45 \\
    \bottomrule
  \end{tabular}
  \end{adjustbox}
  \vspace{11pt}
\end{table}

\begin{figure}[t]
  \centering
  \includegraphics[width=\linewidth]{fig5_tree_vs_alpha.png}
  \caption{Optimal tree depth, maximum level width, and node count versus $\alpha$ (log axis in $1-\alpha$) on both platforms (H2).}
  \label{fig:trees}
\end{figure}

\begin{figure}[t]
  \centering
  \includegraphics[width=\linewidth]{fig4_surface_tesla_t4_float32.png}
  \caption{Predicted speedup over tree size $n$ and depth cap $d$ on T4 FP32 for $\alpha = 0$, the selected $\alpha = 0.995$, and $\alpha = 0.999$; the star marks the optimum. Rising acceptance moves the optimum from a one-token chain to a deep tree.}
  \label{fig:surface}
\end{figure}

\paragraph{Verdict on H2 (depth: supported on both platforms; width: platform-dependent).} Over the feasible $\alpha$ values, optimal depth rises with $\alpha$ on both platforms (Spearman $\rho(\alpha, d) = 0.65$ on each; $0.84$ if the infeasible $\alpha = 1$ is included). The predicted narrowing at the top holds on FP16 (first-level width $4 \to 3$, $\rho(\alpha, \max w) = -0.67$) but not on FP32, where the low-$\alpha$ optimum is already a width-1 chain and the tree widens as it deepens ($\rho(\alpha, \max w) = +0.48$). The shape that higher acceptance buys therefore depends on the hardware.

\subsection{End-to-end speedup: joint vs.\ sequential (Hypothesis H4)}
Because $Q(\alpha)$ drops sharply from $0.0480$ at $\alpha = 0.995$ to $0.0422$ at $\alpha = 0.998$ ($11.3\%$ below $Q(p)$), all three mean-quality tolerances ($\mathrm{tol} \in \{2\%, 5\%, 10\%\}$) select $\alpha^\star = 0.995$ on both platforms. Tables~\ref{tab:e2e} and~\ref{tab:e2e32} summarize end-to-end wall-clock measurements over 100 test articles $\times$ 3 passes.

\begin{table}[h]
  \caption{End-to-end wall-clock performance on 100 CNN/DailyMail test prompts $\times$ 3 passes (T4 FP16). Speedup is computed relative to the autoregressive baseline measured within the same interleaved run.}
  \label{tab:e2e}
  \centering
  \begin{adjustbox}{max width=\linewidth}
  \begin{tabular}{llccccc}
    \toprule
    Method & Target & Topology & Speedup (95\% CI) & ms / tok & Tok / step & Pred.\ (raw $\to$ calib.) \\
    \midrule
    \textbf{JCSD joint (ours)} & $\alpha = 0.995$ & DP ($n=12, d=3$) & \textbf{1.481} $[1.462, 1.500]$ & \textbf{24.66} & \textbf{3.29} & $1.768 \to 1.422$ \\
    Sequential baseline & $\alpha = 0.995$ & SEQUOIA ($n=10, d=2$) & 1.405 $[1.391, 1.418]$ & 26.01 & 2.70 & $1.714 \to 1.387$ \\
    SEQUOIA & Exact $p$ & DP ($n=10, d=2$) & 1.368 $[1.349, 1.387]$ & 26.71 & 2.33 & $1.486 \to 1.234$ \\
    Faster Cascades chain & $\alpha = 0.995$ & Chain ($\gamma = 3$) & 1.352 $[1.330, 1.372]$ & 27.02 & 2.80 & $1.606 \to 1.315$ \\
    Speculative chain ($\gamma=3$) & Exact $p$ & Chain ($n=4, d=3$) & 1.150 $[1.127, 1.174]$ & 31.76 & 2.21 & $1.262 \to 1.075$ \\
    Speculative chain ($\gamma=5$) & Exact $p$ & Chain ($n=6, d=5$) & 0.963 $[0.937, 0.991]$ & 37.95 & 2.45 & $1.053 \to 0.920$ \\
    Speculative chain ($\gamma=7$) & Exact $p$ & Chain ($n=8, d=7$) & 0.814 $[0.793, 0.837]$ & 44.87 & 2.56 & $0.876 \to 0.782$ \\
    Autoregressive (AR) & Exact $p$ & Single token ($n=1$) & 1.000 & 36.53 & 1.00 & $1.000 \to 1.000$ \\
    \bottomrule
  \end{tabular}
  \end{adjustbox}
  \vspace{11pt}
\end{table}

\begin{table}[h]
  \caption{End-to-end wall-clock performance on T4 FP32 (100 prompts $\times$ 3 passes). Rows marked @FP16 run the FP16 platform's optima on FP32 (transfer, H3). Calibration on FP32 fits $c_0 = c_1 = 0$, so predictions are uncalibrated.}
  \label{tab:e2e32}
  \centering
  \begin{adjustbox}{max width=\linewidth}
  \begin{tabular}{llccccc}
    \toprule
    Method & Target & Topology & Speedup (95\% CI) & ms / tok & Tok / step & Predicted \\
    \midrule
    \textbf{JCSD joint (ours)} & $\alpha = 0.995$ & DP ($n=21, d=5$) & \textbf{1.375} $[1.355, 1.398]$ & 36.38 & 4.25 & 1.542 \\
    Sequential baseline & $\alpha = 0.995$ & SEQUOIA ($n=2, d=1$) & 1.185 $[1.179, 1.191]$ & 42.23 & 1.77 & 1.474 \\
    SEQUOIA & Exact $p$ & DP ($n=2, d=1$) & 1.202 $[1.193, 1.212]$ & 41.63 & 1.63 & 1.344 \\
    Faster Cascades chain & $\alpha = 0.995$ & Chain ($\gamma = 3$) & \textbf{1.415} $[1.399, 1.431]$ & \textbf{35.37} & 2.82 & 1.493 \\
    Speculative chain ($\gamma=3$) & Exact $p$ & Chain ($n=4, d=3$) & 1.178 $[1.157, 1.201]$ & 42.46 & 2.22 & 1.173 \\
    Speculative chain ($\gamma=5$) & Exact $p$ & Chain ($n=6, d=5$) & 1.011 $[0.988, 1.034]$ & 49.49 & 2.45 & 1.018 \\
    Speculative chain ($\gamma=7$) & Exact $p$ & Chain ($n=8, d=7$) & 0.777 $[0.758, 0.797]$ & 64.40 & 2.53 & 0.868 \\
    Joint @FP16 & $\alpha = 0.995$ & DP ($n=12, d=3$) & 1.295 $[1.283, 1.306]$ & 38.64 & 3.31 & --- \\
    SEQUOIA @FP16 & Exact $p$ & DP ($n=10, d=2$) & 1.055 $[1.044, 1.067]$ & 47.41 & 2.34 & --- \\
    Autoregressive (AR) & Exact $p$ & Single token ($n=1$) & 1.000 & 50.03 & 1.00 & 1.000 \\
    \bottomrule
  \end{tabular}
  \end{adjustbox}
  \vspace{11pt}
\end{table}

\begin{figure}[t]
  \centering
  \includegraphics[width=0.75\linewidth]{fig6_pareto_tesla_t4_float16.png}
  \caption{Quality versus measured speedup on T4 FP16 (95\% CI on quality). All exact methods share the target distribution and the cascade methods share $\pi_{0.995}$; quality differences between them are within noise.}
  \label{fig:pareto}
\end{figure}

\paragraph{Verdict on H4 (supported in measurement on both platforms).} On FP16, JCSD achieves \textbf{1.481$\times$} speedup ($24.66$\,ms/token), outperforming the sequential baseline ($1.405\times$, $26.01$\,ms/token) by \textbf{5.4\%} (exceeding the predicted $3.2\%$ gain). In paired prompt-level comparisons, sequential minus joint latency is \textbf{$+1.41$\,ms/token with 95\% CI $[+1.02, +1.79]$}, confirming statistically significant acceleration at identical output distribution $\pi_{0.995}$ (all exact methods' paired ROUGE-2 differences vs.\ AR include zero; Figure~\ref{fig:pareto}). The gain stems directly from topology adaptation: SEQUOIA's depth-2 tree caps yield at $2.70$ tokens/step, whereas JCSD's depth-3 tree harvests $3.29$ tokens/step. On FP32 the gap is larger: \textbf{$1.375\times$ vs.\ $1.185\times$ (+16\%)}, paired difference $+5.68$\,ms/token, 95\% CI $[+5.06, +6.29]$, because SEQUOIA's FP32 tree drafts a single token and cannot exploit the higher acceptance at $\alpha = 0.995$ (1.77 vs.\ 4.25 tokens/step). Under a paired-CI feasibility sensitivity check on FP16, $\alpha = 0.999$ ($n=23, d=5$) becomes admissible, raising predicted joint-over-sequential gain to $6.2\%$. Notably, long exact chains ($\gamma \ge 5$) are slower than or equal to AR on both platforms due to the $10$--$13$\,ms per-level draft launch overhead.

\paragraph{Caveat: the predicted joint optimum is not the fastest method on FP32.} The Faster Cascades chain ($\gamma = 3$, $\alpha = 0.995$) measured $1.415\times$, $2.9\%$ faster than the joint optimum (per-run latency difference $+1.08$\,ms/token for joint, 95\% CI $[+0.40, +1.75]$). A chain belongs to the joint search space, so this is a ranking error of the cost model, which had predicted the joint tree $3.3\%$ ahead. Section~\ref{sec:costmodel} traces it to two identifiable modeling gaps. On FP16, the predicted and measured rankings agree on the winner.

\subsection{Cross-platform hardware dependence (Hypothesis H3)}
Section~\ref{sec:profile} and Table~\ref{tab:trees} show that the two platforms select very different trees for the same acceptance statistics. To measure the cost of ignoring this, we ran the FP16 platform's optima on FP32 (Table~\ref{tab:e2e32}, rows @FP16; Figure~\ref{fig:transfer}).

\begin{figure}[t]
  \centering
  \includegraphics[width=0.55\linewidth]{fig8_transfer.png}
  \caption{Transfer matrix of the joint optimum (H3): rows give the platform whose optimum is used, columns the platform it runs on, entries the measured speedup. The FP32 $\to$ FP16 direction was not measured.}
  \label{fig:transfer}
\end{figure}

\paragraph{Verdict on H3 (supported in the measured direction).} Running the FP16 joint tree on FP32 yields $1.295\times$ instead of $1.375\times$ for FP32's own joint tree, a \textbf{5.9\%} transfer loss; transplanting SEQUOIA's FP16 tree yields $1.055\times$ instead of $1.202\times$, a \textbf{12.2\%} loss (all CIs non-overlapping). The cost model predicts losses in both directions ($7.0\%$ for FP16 $\to$ FP32 and $3.4\%$ for FP32 $\to$ FP16); the reverse direction was not measured end to end. Because both platforms are the same GPU model, this isolates the effect of the precision-dependent cost curve; the two runs also differ in host machine and library versions (Table~\ref{tab:setup}).

\section{Cost-model validation and diagnostic insights}
\label{sec:costmodel}

\begin{figure}[t]
  \centering
  \includegraphics[width=0.55\linewidth]{fig7_pred_vs_meas.png}
  \caption{Predicted (uncalibrated) versus measured speedup for every method on both platforms; the dashed line is $y = x$.}
  \label{fig:predmeas}
\end{figure}

\paragraph{Overhead calibration.} On FP16, uncalibrated predictions from Eq.~\eqref{eq:cost} over-predicted speedup by $7.6\%$--$9.7\%$ for exact methods and $18.8\%$--$22.0\%$ for cascade methods (Table~\ref{tab:e2e}, rightmost column; Figure~\ref{fig:predmeas}). Fitting non-negative least squares yields $c_0 = 0$\,ms and $c_1 = 4.7$\,ms per accepted token, reflecting Python tree-walk and host-device synchronization overhead. Calibration reduces in-sample prediction error to within $1.2\%$--$9.8\%$ and preserves the top-ranked method (joint), although SEQUOIA and the Faster Cascades chain swap places; using measured AR loop latency ($36.53$\,ms) instead of profiled AR latency ($33.9$\,ms) in the numerator further tightens errors to $[-2.7\%, +6.5\%]$. On FP32, raw errors were $-0.4\%$ to $+11.8\%$ for exact methods and $+5.5\%$ to $+24.4\%$ for cascade methods, and the fit returns $c_0 = c_1 = 0$ because the residuals have both signs: a constant overhead cannot explain them.

\paragraph{An unpriced draft pass in cascade mode.} Comparing the \emph{same tree} run in exact mode and in cascade mode isolates the cost of the deepest-leaf draft pass described in Section~4.2. That pass fires whenever the walk reaches the deepest level, so its expected cost is $P(\text{reach depth } d)\, t_{\mathrm{draft}}(1)$, which Eq.~\eqref{eq:cost} omits. Table~\ref{tab:deepleaf} shows that this prediction matches the measured extra step time closely on both platforms. The omission penalizes shallow trees under the cascade most (the walk reaches their deepest level most often), which is why the sequential baseline has the largest prediction error.

\begin{table}[h]
  \caption{Extra measured step time of cascade mode over exact mode for the same tree, against the predicted cost of the deepest-leaf draft pass, $P(\text{reach depth } d)\, t_{\mathrm{draft}}(1)$. Residual = measured step time minus the uncalibrated forward-cost model.}
  \label{tab:deepleaf}
  \centering
  \begin{adjustbox}{max width=\linewidth}
  \begin{tabular}{llcccc}
    \toprule
    Platform & Tree & Residual, exact $p$ (ms) & Residual, $\alpha = 0.995$ (ms) & Measured extra (ms) & Predicted extra (ms) \\
    \midrule
    T4 FP16 & Chain $\gamma = 3$ & $+8.5$ & $+13.9$ & $+5.4$ & $5.0$ \\
    T4 FP16 & SEQUOIA tree ($n=10$) & $+8.0$ & $+15.9$ & $+7.9$ & $8.2$ \\
    T4 FP32 & Chain $\gamma = 3$ & $-10.5$ & $-5.2$ & $+5.3$ & $5.3$ \\
    T4 FP32 & SEQUOIA tree ($n=2$) & $+2.5$ & $+9.5$ & $+7.0$ & $8.4$ \\
    \bottomrule
  \end{tabular}
  \end{adjustbox}
  \vspace{11pt}
\end{table}

\paragraph{Profile overstatement for small trees on FP32.} Exact-mode chains with $n = 4$--$6$ ran $10$--$13$\,ms per step faster than the FP32 profile predicts (negative residuals in Table~\ref{tab:deepleaf}), coinciding with the step in $t_{\mathrm{verify}}$ at $n = 4$ (Table~\ref{tab:profile}), whereas the $n = 21$ joint tree matched its prediction within $0.2$\,ms. Together with the unpriced draft pass, this under-predicts the Faster Cascades chain relative to the joint tree and explains the FP32 ranking error in Section~6.4. Both effects are fixable in the cost model (an extra $P(\text{reach})\, t_{\mathrm{draft}}$ term; profiling at the decode-time context length), which we leave to future work.

\paragraph{Testing structural assumptions.} End-to-end visit logs revealed two notable empirical phenomena:
\begin{itemize}
    \item \textbf{Mild violation of positional independence:} In live decoding under exact $p$, first-child acceptance $p_{d,1}$ increases monotonically with tree depth $d$: $0.611$ (root), $0.625$ ($d=1$), $0.646$ ($d=2$), $0.666$ ($d=3$), $0.690$ ($d=4$), $0.699$ ($d=5$), and $0.743$ ($d=6$) on FP16, and $0.616$ to $0.773$ on FP32, versus $0.636$ teacher-forced. Surviving to deep tree levels naturally conditions on high-predictability contexts, whereas higher ranks remain stable ($p_2 \approx 0.090, p_3 \approx 0.044$). The close agreement between platforms also supports reusing the FP16 acceptance statistics on FP32.
    \item \textbf{High sibling waste in static trees:} At branching nodes, the first child was accepted—rendering all drafted siblings unused—\textbf{63\%} of the time for SEQUOIA and \textbf{79\%} (FP16) to \textbf{80\%} (FP32) of the time for JCSD at $\alpha = 0.995$. This strongly motivates confidence-adaptive dynamic tree expansion as future work.
\end{itemize}

\section{Development engineering lessons}

Iterative development surfaced four subtle systems pitfalls worth highlighting:
\begin{enumerate}
    \item \textbf{Logarithmic grid spacing near $\alpha = 1$:} Because quality improvement over raw target sampling persisted up to $\alpha = 0.99$, uniform grids missed the active constraint boundary; spacing $1 - \alpha$ logarithmically ($10^{-2}$ to $10^{-3}$) was essential to resolve the trade-off.
    \item \textbf{Silent artifact collisions:} Formatting $\alpha$ with two decimal places (\texttt{a\%.2f}) caused $\alpha = 0.995$ and $0.998$ to collide with existing files (\texttt{a0.99}, \texttt{a1.00}), causing resumable scripts to skip execution silently until precision-adaptive naming was added.
    \item \textbf{Framework API drift:} Upgrading to Transformers 5.x changed \texttt{DynamicCache.crop()} from taking the number of tokens to keep to taking a negative count of tokens to remove, and renamed GPT-NeoX's output projection to \texttt{lm\_head}, underscoring the necessity of comprehensive $\chi^2$ integration tests.
    \item \textbf{Measurement hygiene on shared notebooks:} An interrupted cell left a benchmark process running on the GPU, which surfaced as an out-of-memory error. Every profile now stores an \texttt{nvidia-smi} snapshot, which let us verify that the FP32 profile ran with no competing process.
\end{enumerate}

\section{Limitations and future work}

Our study is currently bounded by several constraints: (1) evaluation on a single model pair (Pythia-160M $\to$ 2.8B), summarization task (CNN/DailyMail), and temperature ($T=1.0$); (2) batch size 1 inference, whereas larger batch sizes shift the compute-to-memory knee of $t_{\mathrm{verify}}(n)$; (3) Python orchestration overhead ($4.7$\,ms per accepted token on FP16), which a C++/CUDA-graph verifier would eliminate, likely shifting optimal trees even deeper; (4) finite-sample CI width ($\pm 0.01$ ROUGE-2 over 200 prompts), which prevents resolving fine tolerance distinctions ($2\%$ vs.\ $5\%$ vs.\ $10\%$); (5) in-sample overhead calibration; (6) the second platform is the same GPU model at a different precision (P100 GPUs were no longer available on Kaggle), run on a different host with different library versions, and reuses quality and acceptance measured in FP16; (7) the transfer matrix was measured in one direction only (FP16 $\to$ FP32); and (8) the cost model omits the deepest-leaf draft pass, which caused a ranking error on FP32. Future work will add that term to the cost model, measure the reverse transfer direction, extend evaluation to Qwen2.5 (0.5B $\to$ 3B) on GSM8K at $T=0.6$, and replace static DP trees with dynamic, draft-entropy-gated branching to reclaim the $79$--$80\%$ sibling waste observed under high-$\alpha$ cascades.

\section{Conclusion}

We presented JCSD, a hardware-aware framework that jointly optimizes cascade deferral $\alpha$ and token-tree topology $T$ under a quality budget. By decoupling hardware-independent quality and acceptance profiling from hardware-priced dynamic programming, JCSD makes joint optimization computationally tractable. On Pythia-160M $\to$ Pythia-2.8B (CNN/DailyMail, Tesla T4 FP16), JCSD adapts to the elevated acceptance rate at $\alpha = 0.995$ by transitioning from a depth-2 tree to a depth-3 tree, achieving a measured \textbf{1.481$\times$} speedup over autoregressive decoding and outperforming sequential tuning by \textbf{5.4\%} ($1.41$\,ms/token reduction) at identical output quality. On the same GPU in FP32, the optimal trees change sharply, transplanting the FP16 optimum loses $5.9\%$, and joint search beats sequential by $16\%$, confirming that the best $(\alpha, T)$ depends on the hardware. The FP32 results also show the method's main limitation: the joint optimum is only as good as its cost model, which missed a cheaper chain by $2.9\%$ because of a draft pass it does not price.

\section*{Author contributions}

All five authors contributed equally to this project. The individual contributions are:

\begin{itemize}
    \item \textbf{Ayanna Mondal:} Surveyed the literature on speculative decoding, SEQUOIA and Faster Cascades, and framed the gap between them. Implemented the TokenV3 cascade target $\pi_\alpha$ and the batched float64 verifier (\texttt{AcceptCounter}), and wrote the $\chi^2$ tests that check exact sampling.
    \item \textbf{Swastik Pramanik:} Formulated the constrained joint $(\alpha, T)$ objective and the decoupling argument. Implemented the expected-token dynamic program, the tree search (\texttt{run\_search.py}), the analytical cost model, and the joint versus sequential selection.
    \item \textbf{Takshay Bansal:} Built the shared \texttt{TreeDecoder}, the tree attention masks, and the KV-cache surgery (\texttt{keep\_path}, \texttt{crop\_to}). Ported the code to Transformers 5.x, maintained the test suite, and ran the full pipeline on Google Colab (FP16) and Kaggle (FP32).
    \item \textbf{Youkta Mandavkar:} Implemented the quality and acceptance measurement (\texttt{run\_quality.py}, ROUGE-2) and the interleaved end-to-end benchmark (\texttt{run\_e2e.py}) with paired bootstrap confidence intervals. Analysed hypotheses H1, H2 and H4.
    \item \textbf{Shirshendu Sengupta:} Profiled the hardware latencies on FP16 and FP32 (\texttt{run\_profile.py}), fitted the overhead calibration (\texttt{calibrate.py}), and diagnosed the unpriced deepest-leaf draft pass. Analysed cross-platform transfer (H3) and produced the figures and tables.
\end{itemize}

All authors took part in the design discussions, reviewed each other's code and results, and co-wrote this report and the presentation.

% \begin{ack}
% This work was completed as Project 3 for the Hardware-Aware Deep Learning (HDL) course at the Indian Institute of Technology Guwahati. Compute resources were provided via Google Colaboratory and Kaggle.
% \end{ack}


\medskip
{
\small
\begin{thebibliography}{9}

\bibitem[Biderman et~al.(2023)]{biderman2023pythia}
Biderman, S., Schoelkopf, H., Anthony, Q.~G., Bradley, H., O'Brien, K., Hallahan, E., Khan, M.~A., Purohit, S., Prashanth, U.~S., Raff, E., et~al. (2023).
\newblock Pythia: A suite for analyzing large language models across training and scaling.
\newblock In \emph{International Conference on Machine Learning (ICML)}, pp.\ 2397--2430.

\bibitem[Chen et~al.(2023)]{chen2023accelerating}
Chen, C., Borgeaud, S., Irving, G., Lespiau, J.-B., Sifre, L., \& Jumper, J. (2023).
\newblock Accelerating large language model decoding with speculative sampling.
\newblock \emph{arXiv preprint arXiv:2302.01318}.

\bibitem[Chen et~al.(2024)]{chen2024sequoia}
Chen, Z., May, A., Svirschevski, R., Huang, Y., Ryabinin, M., Jia, Z., \& Chen, B. (2024).
\newblock SEQUOIA: Scalable, robust, and hardware-aware speculative decoding.
\newblock In \emph{Advances in Neural Information Processing Systems (NeurIPS)}, vol.\ 37.

\bibitem[Leviathan et~al.(2023)]{leviathan2023fast}
Leviathan, Y., Kalman, M., \& Matias, Y. (2023).
\newblock Fast inference from transformers via speculative decoding.
\newblock In \emph{International Conference on Machine Learning (ICML)}, pp.\ 19274--19286.

\bibitem[Lin(2004)]{lin2004rouge}
Lin, C.-Y. (2004).
\newblock ROUGE: A package for automatic evaluation of summaries.
\newblock In \emph{Text Summarization Branches Out}, pp.\ 74--81.

\bibitem[Narasimhan et~al.(2025)]{narasimhan2025faster}
Narasimhan, H., Jitkrittum, W., Rawat, A.~S., Kim, S., Gupta, N., Menon, A.~K., \& Kumar, S. (2025).
\newblock Faster cascades via speculative decoding.
\newblock In \emph{International Conference on Learning Representations (ICLR)}.

\bibitem[See et~al.(2017)]{see2017get}
See, A., Liu, P.~J., \& Manning, C.~D. (2017).
\newblock Get to the point: Summarization with pointer-generator networks.
\newblock In \emph{Proceedings of the 55th Annual Meeting of the Association for Computational Linguistics (ACL)}, pp.\ 1073--1083.

\end{thebibliography}
}

%%%%%%%%%%%%%%%%%%%%%%%%%%%%%%%%%%%%%%%%%%%%%%%%%%%%%%%%%%%%
\newpage
\section*{NeurIPS Paper Checklist}

\begin{enumerate}

\item {\bf Claims}
    \item[] Question: Do the main claims made in the abstract and introduction accurately reflect the paper's contributions and scope?
    \item[] Answer: \answerYes{} % Replace by \answerYes{}, \answerNo{}, or \answerNA{}.
    \item[] Justification: The abstract and Section 1 explicitly state the scope (Pythia-160M $\to$ Pythia-2.8B on CNN/DailyMail, Tesla T4 in FP16 and FP32, batch size 1), match the measured speedups reported in Section 6.4 ($1.481\times$ vs.\ $1.405\times$ on FP16; $1.375\times$ vs.\ $1.185\times$ on FP32) and the transfer loss in Section 6.5, and state the FP32 case where a Faster Cascades chain beat the predicted joint optimum.

\item {\bf Limitations}
    \item[] Question: Does the paper discuss the limitations of the work performed by the authors?
    \item[] Answer: \answerYes{}
    \item[] Justification: Section 9 provides a dedicated discussion of eight concrete limitations: single model pair/task/temperature, batch size 1 regime, Python orchestration overhead, finite-sample confidence interval width over 200 calibration prompts, in-sample overhead calibration, a second platform that differs in precision (same GPU model), one-direction transfer measurement, and the unpriced deepest-leaf draft pass in the cost model.

\item {\bf Theory assumptions and proofs}
    \item[] Question: For each theoretical result, does the paper provide the full set of assumptions and a complete (and correct) proof?
    \item[] Answer: \answerYes{}
    \item[] Justification: Section 3 states the TokenV3 target distribution, positional acceptance independence assumption, and exact residual sampling invariant, and Section 7 explicitly tests and quantifies violations of the positional independence assumption in live decoding.

\item {\bf Experimental result reproducibility}
    \item[] Question: Does the paper fully disclose all the information needed to reproduce the main experimental results of the paper to the extent that it affects the main claims and/or conclusions of the paper (regardless of whether the code and data are provided or not)?
    \item[] Answer: \answerYes{}
    \item[] Justification: Section 5 (Table 2) specifies exact Hugging Face model identifiers, prompt truncation rules, generation lengths, dataset splits, $\alpha$ grids, DP caps, and software/GPU versions for both platforms.

\item {\bf Open access to data and code}
    \item[] Question: Does the paper provide open access to the data and code, with sufficient instructions to faithfully reproduce the main experimental results, as described in supplemental material?
    \item[] Answer: \answerYes{}
    \item[] Justification: The complete pipeline, unit tests, and reproduction scripts are publicly hosted in the GitHub repository linked in the author footnote (\texttt{complete-pipeline} branch), using standard public Hugging Face datasets (CNN/DailyMail).

\item {\bf Experimental setting/details}
    \item[] Question: Does the paper specify all the training and test details (e.g., data splits, hyperparameters, how they were chosen, type of optimizer) necessary to understand the results?
    \item[] Answer: \answerYes{}
    \item[] Justification: Section 5 and Table 2 detail the disjoint splits (200 validation prompts for calibration, 100 test prompts for end-to-end evaluation, 2 train prompts for warm-up), temperature, DP bounds, and interleaved evaluation protocol.

\item {\bf Experiment statistical significance}
    \item[] Question: Does the paper report error bars suitably and correctly defined or other appropriate information about the statistical significance of the experiments?
    \item[] Answer: \answerYes{}
    \item[] Justification: Tables 4, 6 and 7 and Section 6.4 report 95\% bootstrap confidence intervals across 100 test prompts $\times$ 3 passes using common random numbers, including the paired latency differences ($+1.41$\,ms/token, 95\% CI $[+1.02, +1.79]$ on FP16; $+5.68$\,ms/token, 95\% CI $[+5.06, +6.29]$ on FP32).

\item {\bf Experiments compute resources}
    \item[] Question: For each experiment, does the paper provide sufficient information on the computer resources (type of compute workers, memory, time of execution) needed to reproduce the experiments?
    \item[] Answer: \answerYes{}
    \item[] Justification: Section 5 and Section 6 specify the GPU hardware (NVIDIA Tesla T4 16\,GB on Google Colab and on Kaggle), precision (FP16 and FP32), profiling repetitions (30), and runtime characteristics ($1.3$ and $2.3$ decoding GPU-hours for the FP16 and FP32 end-to-end benchmarks).

\item {\bf Code of ethics}
    \item[] Question: Does the research conducted in the paper conform, in every respect, with the NeurIPS Code of Ethics \url{https://neurips.cc/public/EthicsGuidelines}?
    \item[] Answer: \answerYes{}
    \item[] Justification: This work focuses strictly on inference efficiency and latency optimization for existing open-weight language models on standard public summarization benchmarks and conforms to the NeurIPS Code of Ethics.

\item {\bf Broader impacts}
    \item[] Question: Does the paper discuss both potential positive societal impacts and negative societal impacts of the work performed?
    \item[] Answer: \answerNA{}
    \item[] Justification: This is foundational systems research improving LLM decoding throughput and energy efficiency without introducing new generative capabilities or deployment risks beyond those of the underlying base models.

\item {\bf Safeguards}
    \item[] Question: Does the paper describe safeguards that have been put in place for responsible release of data or models that have a high risk for misuse (e.g., pre-trained language models, image generators, or scraped datasets)?
    \item[] Answer: \answerNA{}
    \item[] Justification: No new pre-trained weights or scraped datasets are released; the work evaluates existing public Pythia models on CNN/DailyMail.

\item {\bf Licenses for existing assets}
    \item[] Question: Are the creators or original owners of assets (e.g., code, data, models), used in the paper, properly credited and are the license and terms of use explicitly mentioned and properly respected?
    \item[] Answer: \answerYes{}
    \item[] Justification: All existing models (EleutherAI Pythia, Apache-2.0) and datasets (CNN/DailyMail, Apache-2.0) are properly cited in the References section and used in accordance with their open licenses.

\item {\bf New assets}
    \item[] Question: Are new assets introduced in the paper well documented and is the documentation provided alongside the assets?
    \item[] Answer: \answerYes{}
    \item[] Justification: The open-source JCSD codebase includes documented pipeline modules, 84 automated tests, version-bounded dependency requirements, and the exact library versions recorded in every result artifact.

\item {\bf Crowdsourcing and research with human subjects}
    \item[] Question: For crowdsourcing experiments and research with human subjects, does the paper include the full text of instructions given to participants and screenshots, if applicable, as well as details about compensation (if any)?
    \item[] Answer: \answerNA{}
    \item[] Justification: The paper does not involve crowdsourcing or human subjects.

\item {\bf Institutional review board (IRB) approvals or equivalent for research with human subjects}
    \item[] Question: Does the paper describe potential risks incurred by study participants, whether such risks were disclosed to the subjects, and whether Institutional Review Board (IRB) approvals (or an equivalent approval/review based on the requirements of your country or institution) were obtained?
    \item[] Answer: \answerNA{}
    \item[] Justification: The paper does not involve human subjects or require IRB approval.

\item {\bf Declaration of LLM usage}
    \item[] Question: Does the paper describe the usage of LLMs if it is an important, original, or non-standard component of the core methods in this research?
    \item[] Answer: \answerNA{}
    \item[] Justification: Large language models (Pythia-160M and Pythia-2.8B) are the subject of inference acceleration rather than an external tool used to generate the research methodology.

\end{enumerate}

\end{document}
